\section{Limitations}\label{sec:limitations}

The following limitations are stated in full; items 16 to 20 concern the extension, and item 21 the distinction between a calibrated and a useful interval.

\begin{enumerate}
\item \textbf{Procedural boundary.} The hold-out countries are untouched with respect to countries only, and they share calendar months with the design sample. The methods, the form of the contrasts, the horizons, the margin of 2~pp and the thresholds of 3~pp and 7~pp were chosen on the 13 design countries, and the author had seen the post-2025-06 months of these countries before the amendments. The rules R1 to R5 were written after the exploratory results at both warm-ups were known. The pre-specification is a repository commit history, kept by one author and self-attested, with all entries on 2026-11-18 or 2026-11-19, and not an external registry or an independent timestamp.
\item \textbf{Power and the ceiling effect.} The minimum detectable difference is undefined on the full first hold-out sample for all four primary contrasts. All null outcomes are therefore inconclusive and are not evidence of equality. The 2~pp equivalence margin is about the size of ACI's own calibration error, so the equivalence verdict for DRACP against ACI is weak by construction and the contrasts against ACI cannot be powered; null outcomes were likely before the run. The MDE simulation has not been reviewed independently, and it assumes that the observed hit pattern of the baseline describes the true dependence and that dependence beyond the block length is zero.
\item \textbf{Small panel.} The first hold-out panel has 7 countries; Ghana and Angola have high-inflation episodes, so pooled results can be driven by few units. S1 shows that removing one country changes the sign of ACI against rolling at $h=6$ on the hold-out panel in one case of seven.
\item \textbf{One point-forecaster family dominates the decision rules (Sections~\ref{sec:explor} and~\ref{sec:holdout}).} LightGBM is the base model for the primary family of R1--R5. ARIMA and naive B give the same direction in S2, but they were not part of the pre-specified decision rules. In the extension the four base models share the same countries and months, so they are not independent replications, and the Holm correction is applied within each base model and not across them.
\item \textbf{Data revisions and vintages were not checked.} The IMF publishes dated CPI vintages, which allow a real-time revision check, but this was not done. Evaluation uses the revised data of the snapshot date. Small revisions for Russia and Mongolia are noted between the two snapshots, but their effect on intervals was not studied. Nigeria's series ends at 2024-12 because of a methodology break, so that country has no recent months.
\item \textbf{No systematic literature search was run.} The statements about novelty and the state of the literature are based on the author's reading only, and the bibliography is limited to a fixed list of works, several of whose entries are marked as needing verification.
\item \textbf{``Commodity-dependent'' is not defined by a measurable criterion.} No export-share rule was applied to the 13 design countries or to the hold-out candidates; Georgia and Uruguay are borderline. The hold-out sets were determined by IMF data availability under fixed written rules, and 7 of 14 candidates qualified for the first panel.
\item \textbf{Reproduction of DRACP is partial.} Five countries with annual HICP inflation, one seed, and a mean coverage of ACI 0.874 and DRACP 0.871 (against 0.884 and 0.890 reported over 48 series with 20 seeds). The criterion of reproducing a published number within 2~pp could not be tested because the published numbers are averages over different series. The DRACP configuration in the benchmark follows the reproduction script and was not tuned. The description of DRACP in Section~\ref{sec:related} is as described by its authors and was read, not re-derived.
\item \textbf{Calibration error is noisy and its bootstrap is not validated at the boundary.} CE is biased upward in small samples, and percentile intervals are not symmetric around the estimate. The validity of the CE bootstrap at the boundary of the absolute value (where coverage is close to 0.90) was not checked by simulation, so the size of the test is unknown. The choice of the 2~pp equivalence margin and of the 3~pp threshold of R4 is not justified beyond being fixed in advance.
\item \textbf{The stress analysis is secondary and exploratory.} The exchange-rate flag is a weak proxy of inflation shocks (lift below 2), is affected by exchange-rate regimes, and many deciles, methods and horizons were examined. Pegged and managed currencies were not analysed separately.
\item \textbf{Items in the protocol that are not reported.} The central-bank target baseline, the exchange-rate-regime split, and per-country heterogeneity beyond coverage ranges are mentioned in the protocol; the results documents do not report them. The foundation-model forecaster (Chronos) is reported as an extension, with the limitations of items 16 to 20. The native ARIMA and ETS intervals are reported only as a post hoc secondary analysis, and Kazakhstan only descriptively.
\item \textbf{The post-2025-06 months of the 13 countries are degenerate} (at most 16 months per unit, blocks of 12) and are not used for any claim. The $h=12$ horizon has few non-overlapping windows (80 non-overlapping cells on the first hold-out panel in run A).
\item \textbf{The platform is not the benchmark.} The platform uses an exponentially weighted moving-average point forecast and a smaller minimum calibration size, so the benchmark results apply to its interval method but not to its point forecast. The module is an engineering application and not a scientific contribution.
\item \textbf{Point forecasts matter for interval width but were not optimised.} All hyperparameters were fixed in advance, and the point models are simple; better point forecasts could change absolute calibration and width.
\item \textbf{The post hoc additions have no multiplicity correction.} S6 (Kazakhstan) is descriptive with about 137 origins of one design-panel country and no test; S7 (native intervals) uses unadjusted $p$-values and one base model for the contrast.
\item \textbf{The ACI-against-rolling result depends on the panel, and the reason is not known.} R2 is met at warm-up 48 on the second hold-out panel (at $h=3$ for LightGBM, ARIMA and naive B, Holm 0.008, 0.004, 0.004, and at $h=6$ for ARIMA and naive B, 0.030 and 0.039; Chronos at $h=3$, Holm 0.036, has an undefined MDE and is inconclusive under E5) and is not shown (MDE undefined, inconclusive) on the design panel and the first hold-out panel, and not significant (MDE undefined in every cell, inconclusive) on the earlier 20 countries re-run from 2010-01. Candidate explanations (a shorter series, more volatile and lower-income economies, a warm-up that takes a larger share of the series, a rolling window with a larger calibration error) were not tested. The panels are not pooled, and no panel is selected.
\item \textbf{The second hold-out panel is short.} It has 198 months from 2010-01, so 89 and 83 pooled months at $h=3$ and $h=6$ at warm-up 48, and 41 and 35 at warm-up 96. Warm-up 96 shows nothing on it, and the minimum detectable difference is undefined for every contrast at warm-up 96 and for every DRACP minus ACI contrast at warm-up 48. Several of its countries are candidates that failed the first hold-out rule on history length, so the two panels differ in the length of their series by construction.
\item \textbf{Chronos: possible leakage and a thin clean sample.} Chronos may have seen public CPI series in pretraining. The pretraining cutoff was not verified, and 2025-01 is a conservative assumption. Accuracy claims rest on 22 calendar months (567 country-month cells per horizon), with a short block length. Calibration contrasts for Chronos use all months with the flag, and no analysis restricted to clean months is reported. Only the smallest Chronos-Bolt model was used, zero-shot; larger models, fine-tuning and covariates were not tried. On the earlier 20 countries the four base models were run, but calibration contrasts were analysed only for LightGBM and Chronos, which narrows the E-B analysis of the protocol.
\item \textbf{The language-model explainer evaluation is narrow.} One model, one prompt template, temperature 0, one run, 13 countries; the first 37 of the 936 answers used different token budgets. The pre-specified metrics had two defects found after the results were seen; the corrected metrics are post hoc, were checked on 60 answers by a single reader that was the model itself and knew the flags, and still produce false positives (list items under a negating lead-in). The 100\% grounded faithfulness is a ceiling of what the metrics see, and no claim is made that grounding works in general or that the model is faithful in general. The ungrounded condition is a floor (omission 100\% by construction). The human review of fluency and terminology (90 answers) was not done, and Kazakh fluency needs a native speaker.
\item \textbf{Extension specified with knowledge of earlier results.} The eighth entry and the rules E1--E5 were written after R1--R5 had been applied and all earlier results, including R2 on the first hold-out panel, were known; the extension is a replication on new countries in the sense of Section~\ref{sec:boundary} and not a blind test.
\item \textbf{Calibrated does not mean useful.} Coverage close to the nominal level says nothing about whether the interval is narrow enough for a decision. Any method can reach 90 percent with very wide intervals, which is why width and the interval score are reported next to CE (Section~\ref{sec:width}). The reported widths show the cost: the pooled mean width of ACI intervals is 4.28 percentage points of cumulative growth at $h=3$ and 15.5 at $h=12$ in the exploratory data, and the better calibration of the adaptive methods is paid for in width or interval score at long horizons. At $h=12$ an interval with honest 90 percent coverage may be too wide for practical policy or business use. The thesis does not establish a width at which an interval becomes useful, and it does not compare widths with a decision-relevant threshold.
\end{enumerate}
